\item On étudie le lien entre la pression artérielle $Y$ et l'âge $x_1$. Dans un
premier modèle, on obtient
\[
  \widehat{Y}=92+0{,}72x_1.
\]
Après avoir ajouté l'indice de masse corporelle $x_2$, on obtient
\[
  \widehat{Y}=78+0{,}41x_1+1{,}15x_2.
\]

\begin{enumerate}[label=\alph*)]
  \item Interprétez le coefficient associé à l'âge dans chacun des modèles.
  \item Pourquoi les deux coefficients associés à l'âge peuvent-ils être
        différents?
  \item Peut-on déterminer, à partir de ces seules équations, lequel des deux
        modèles décrit le mieux les données? Expliquez.
\end{enumerate}

\prof{
\begin{enumerate}[label=\alph*)]
  \item Dans le premier modèle, une année d'âge supplémentaire est associée à
        une augmentation moyenne de $0{,}72$ unité de pression artérielle. Il
        s'agit d'une association non ajustée pour l'indice de masse corporelle.

        Dans le second modèle, à indice de masse corporelle constant, une année
        d'âge supplémentaire est associée à une augmentation moyenne de
        $0{,}41$ unité de pression artérielle.

  \item L'âge et l'indice de masse corporelle peuvent être associés entre eux,
        et chacun peut être associé à la pression artérielle. Dans le premier
        modèle, le coefficient de l'âge peut alors capter une partie de
        l'association attribuable à l'indice de masse corporelle. Dans le
        second modèle, l'effet de $x_1$ est estimé en maintenant $x_2$ constant.

  \item Non. Les équations fournissent les coefficients estimés, mais pas la
        qualité de l'ajustement. On pourrait notamment examiner les résidus et
        des mesures comme la somme des carrés résiduels ou le coefficient de
        détermination ajusté. Le choix doit aussi tenir compte de l'objectif de
        l'analyse et des postulats du modèle. (On y reviendra dans les prochains chapitres).
\end{enumerate}
}
